\section*{अध्याय \ref{ch:g9:earth-resources} --- पृथ्वी से मिलने वाले संसाधन}

\begin{solution}{exo:g9:earth-resources:1}
(क) \omterm{def:g9:earth-resources:aggregate}{समुच्चय}। (ख) \omterm{def:g9:earth-resources:reserve}{भंडार}। (ग) \omterm{def:g9:earth-resources:ore-deposit}{अयस्क निक्षेप}। (घ) \omterm{def:g9:earth-resources:ore-deposit}{सांद्रण गुणक}।
\end{solution}

\begin{solution}{exo:g9:earth-resources:2}
(क) \omterm{def:g1:landscapes-around-us:water}{नदी} की \omterm{def:g3:stones-pebbles-sand:pebble}{बजरी} में। (ख) दरारों में, गरम पानी से। (ग) \omterm{def:g7:volcanoes:magma}{मैग्मा} में।
(घ) यह \omterm{def:g1:landscapes-around-us:water}{नदी} की \omterm{def:g3:stones-pebbles-sand:pebble}{बजरी} वाले निक्षेप का पहला चरण है: मुक्त हुआ सोना फिर \omterm{def:g1:landscapes-around-us:water}{नदी} द्वारा बहाया
जाता है।
\end{solution}

\begin{solution}{exo:g9:earth-resources:3}
(ख), (घ), (क), (ङ), (ग)।
\end{solution}

\begin{solution}{exo:g9:earth-resources:4}
(क) नवीकरणीय नहीं: उसे बनने में लाखों वर्ष लगते हैं। (ख) नवीकरणीय: जो लिया जाता है उसे
\omterm{def:g4:clouds-rain-snow:precipitation}{बारिश} फिर से भर देती है। (ग) नवीकरणीय नहीं: \omterm{def:g9:earth-resources:ore-deposit}{अयस्क निक्षेप} बहुत लंबे समयों में बनता है।
(घ) उस दर पर नवीकरणीय नहीं: \omterm{def:g6:water-underground:water-table}{भौम जलस्तर} गिरता है और \omterm{def:g6:water-underground:spring}{कुएँ} आख़िर में सूख जाते
हैं।
\end{solution}

\begin{solution}{exo:g9:earth-resources:5}
(क) $4 \div 0.007 \approx 571$, लगभग $10^3$ ($10^2$ और $10^3$ के बीच,
$10^3$ के ज़्यादा क़रीब)। (ख) $0.5 \div 0.0002 = 2500$, लगभग $10^3$। (ग) $20
\div 5.6 \approx 3.6$, लगभग 4: $10^0$ और $10^1$ के बीच।
(घ) लोहा \omterm{def:g6:inside-the-earth:layers}{भूपर्पटी} में पहले से ही बहुत प्रचुर है: साधारण \omterm{def:g3:stones-pebbles-sand:rock}{चट्टान} से केवल कुछ गुना ज़्यादा समृद्ध
अयस्क काफ़ी है।
\end{solution}

\begin{solution}{exo:g9:earth-resources:6}
(क) $19.3 \div 2.5 \approx 7.7$ गुना। (ख) सोना: $19.3 \times 0.01 =
\qty{0.193}{g}$; \omterm{def:g3:stones-pebbles-sand:rock}{चट्टान}: $2.5 \times 0.01 = \qty{0.025}{g}$। (ग) जब पूरा तल
खिसकता है, तो सोने के भारी कण हल्की \omterm{def:g3:stones-pebbles-sand:sand}{रेत} और \omterm{def:g3:stones-pebbles-sand:pebble}{कंकड़ों} के बीच से नीचे डूबते हैं, ठोस आधार-शैल
तक, जहाँ वे रुक जाते हैं।
\end{solution}

\begin{solution}{exo:g9:earth-resources:7}
(क) सोना, टिन और क्रोमियम। (ख) $10^1$ और $10^2$ के बीच, $10^2$ के
क़रीब। (ग) $3000 \div 3.7 \approx 800$ गुना। (घ) क्रोमियम (या टिन, या सोना):
\omterm{def:g9:earth-resources:ore-deposit}{सांद्रण गुणक} जितना ऊँचा, उतनी ही दुर्लभ वे जगहें जहाँ प्रकृति ने धातु को काफ़ी सांद्रित
किया है।
\end{solution}

\begin{solution}{exo:g9:earth-resources:8}
(क) $18 \times 365 = \qty{6570}{kg}$, लगभग \qty{6.6}{t}। (ख) $5 \times
10^{10}$ टन। (ग) $25000 \div 18 \approx 1389$ दिन, लगभग चार वर्ष।
(घ) \omterm{def:g1:landscapes-around-us:water}{नदी} से खोदी गई \omterm{def:g3:stones-pebbles-sand:sand}{रेत} अब नीचे की ओर \omterm{def:g1:landscapes-around-us:water}{तट} तक नहीं जाती; लहरें \omterm{def:g2:water-on-earth:ocean}{समुद्र} \omterm{def:g1:landscapes-around-us:water}{तट} से \omterm{def:g3:stones-pebbles-sand:sand}{रेत} बहाती
रहती हैं, पर उसकी जगह लेने के लिए कम \omterm{def:g3:stones-pebbles-sand:sand}{रेत} पहुँचती
है।
\end{solution}

\begin{solution}{exo:g9:earth-resources:9}
(क) \omterm{def:g2:water-on-earth:ocean}{समुद्र} के नन्हे पौधों और जंतुओं से, जैसे डायटम, जो समुद्री तल के \omterm{def:g3:stones-pebbles-sand:sand}{कीचड़} में दब गए।
(ख) लगभग 60 और \qty{120}{\degreeCelsius} के बीच।
(ग) वे पानी से हल्के होते हैं, जो \omterm{def:g3:stones-pebbles-sand:rock}{चट्टानों} के छिद्रों को भरता है।
(घ) एक \omterm{def:g6:water-underground:permeability}{अपारगम्य} परत, जैसे शेल; वे उसके नीचे इकट्ठा होते हैं, अक्सर किसी उभार के शिखर पर,
तेल के ऊपर गैस।
\end{solution}

\begin{solution}{exo:g9:earth-resources:10}
(क) $2.9 \times 10^7$ टन। (ख) $2.9 \times 10^7 \div 100 = 2.9 \times
10^5$ टन प्रति वर्ष, मोटे तौर पर। (ग) अयस्क अपने अंदर के ताँबे से 100 गुना भारी है:
$2.9 \times 10^7 \times 100 = 2.9 \times 10^9$ टन।
(घ) हर टन अयस्क से केवल \qty{10}{kg} ताँबा मिलता है, इसलिए अरबों टन \omterm{def:g3:stones-pebbles-sand:rock}{चट्टान} खोदनी
पड़ती है।
\end{solution}

\begin{solution}{exo:g9:earth-resources:11}
(क) 23~मिलियन टन प्रति वर्ष पर वह लगभग $280 \div 23
\approx 12$ वर्ष चलता,
1980 के दशक की शुरुआत तक; फिर भी 697~मिलियन टन खोदा गया, उस \omterm{def:g9:earth-resources:reserve}{भंडार} से कहीं ज़्यादा।
(ख) $980 \div 280 = 3.5$ गुना। (ग) $980 \div
23 \approx 43$ वर्ष। (घ) उदाहरण के लिए: नए निक्षेप खोजे जाते हैं; बेहतर तकनीकें या ऊँची
क़ीमतें ज़्यादा कमज़ोर निक्षेपों को खनन लायक़ बना देती हैं।
\end{solution}

\begin{solution}{exo:g9:earth-resources:12}
(क) $5 \times 10^1$ वर्ष। (ख) $2 \times 10^3 \div 50 = 40$ गुना। (ग) $3
\times 10^8 \div (5 \times 10^1) = 6 \times 10^6$ गुना। (घ) कोयले को बनने में उससे
लाखों गुना ज़्यादा समय लगा जितना हमें उसे इस्तेमाल करने में लगता है: प्रकृति उसे हमारे
समय-पैमाने पर फिर से नहीं बनाती।
\end{solution}

\begin{solution}{pb:g9:earth-resources:1}
\textbf{1.} $100 \times 1000 = \num{100000}$~kg $= 1 \times 10^5$~kg।

\textbf{2.} $2000 \times 300 = \num{600000}$~t $= 6 \times 10^5$~t।

\textbf{3.} $6 \times 10^5 \div 10 = 6 \times 10^4$~t प्रति वर्ष।

\textbf{4.} $0.40 \times 6 \times 10^4 = 2.4 \times 10^4$~t \omterm{def:g3:stones-pebbles-sand:sand}{रेत} प्रति
वर्ष।

\textbf{5.} $1200 \times 500 \times 5 = 3 \times 10^6$~m$^3$।

\textbf{6.} उसका दो तिहाई: $2 \times 10^6$~m$^3$।

\textbf{7.} $2 \times 10^6 \times 1.6 = 3.2 \times 10^6$~t।

\textbf{8.} $\num{24000} + \num{136000} = \num{160000}$~t $= 1.6 \times
10^5$~t प्रति वर्ष।

\textbf{9.} $\num{160000} \div 2000 = 80$ गुना ज़्यादा।

\textbf{10.} नहीं: गड्ढा उससे 80 गुना तेज़ी से ख़ाली होता है जितनी तेज़ी से \omterm{def:g1:landscapes-around-us:water}{नदी} \omterm{def:g3:stones-pebbles-sand:sand}{रेत} को फिर से
लाती है। नगर के पैमाने पर यह \omterm{def:g3:stones-pebbles-sand:sand}{रेत} एक सीमित संसाधन है।

\textbf{11.} $\num{24000} \div 4 = \num{6000}$~t \omterm{def:g3:stones-pebbles-sand:sand}{रेत} प्रति वर्ष।

\textbf{12.} $3.2 \times 10^6 \div (1.6 \times 10^5) = 20$: \textbf{20 वर्ष}।
\end{solution}
